
\section{Non-finitariness of FGLS Monad}

\begin{theorem} The Full-Ground Local State Monad is not finitary. \end{theorem}
\begin{proof}
We will show that there exists a filtered category $\mcatC$ and functor $\vF : \mcatC \to \funC\worldC\setC$ such that
\[ \monFLS(\colimit\vF) \not\iso \colimit(\monFLS\vF) \]
in the category of cocones for the functor $\vF$.

Take $\mcatC = \omega$, the posetal category where objects are natural numbers and morphisms are orderings $\vm \leq \vn$.
This category is $\omega$-filtered.
Further define $\vF : \omega \to \funC\worldC\setC$ to be $\vF\vn = \Delta_\vn$ that maps each natural number $\vn$ to the constant functor mapping to the downset of $\vn$ and $\vm \leq \vn$ are mapped to inclusions.

Notice that we can simplify $\colimit\vF$:
\[ (\colimit\vF)\vw = \tup{\quotient{\disunion{\vn \in \nat}\downset\vn}\sim}{\tlistr{\vi \in \downset\vn \mapsto \quot\vn(\vi)}{\vn \in \nat}} 
                    \iso \tup\nat{\tlistr{\vi \in \downset\vn \mapsto \vi}{\vn \in \nat}} \]
which extends to an isomorphism of cocones $\monFLS(\colimit\vF) \iso \tup{\monFLS(\Delta_\nat)}{\tlistr{\monFLS(\Delta_\vn \inclto \Delta_\nat)}{\vn \in \nat}}$.

We also know that:
\[ \colimit(\monFLS\vF)\vw \iso \tup{\quotient{\disunion{\vn \in \nat} \monFLS{\Delta_\vn}\vw}\sim}
                                    {\tlistr{\vf \in \monFLS\Delta_\vn\vw \mapsto \quot\vn(\vf)}{\vn \in \nat}} \]
where the quotienting says that $\tup\vn{\vf \in \monFLS(\Delta_\vn)\vw} \sim \tup\vm{\vg \in \monFLS(\Delta_\vm)\vw}$ iff they are the same morphisms after we compose them with their corresponding inclusions to $\monFLS(\Delta_{\vn + \vm})$.

The above observations boil down to showing that for some $\vw \in \worldC$:
\[ \quotient{\disunion{\vn \in \nat} \monFLS{\Delta_\vn}\vw}\sim \not\iso \monFLS(\Delta_\nat)\vw \]
where the mapping is given by $\quot\vn(\vf \in \monFLS(\Delta_\vn)\vw) \mapsto \monFLS(\Delta_\vn \inclto \Delta_\nat)_\vw(\vf)$.
We show that this mapping is not surjective.

\paragraph{Counterexample}

W.l.g. assume we only have one sort $\phi(\vC) = \tyUnit \tySum \tyRef\vC$ and let $\vw = \va : \vC$ be a singleton world. 
We define $\countFun : \monFLS(\Delta_\nat)\vw$ such that there cannot be element in $\monFLS(\Delta_\vn)\vw$ for some $\vn$ that represents it after post-composition with inclusion.

\[\begin{array}{ll}
    &\countFun : \End{\va : \vC \injto \vw' \in \commaC{\va : \vC}\vU} \heapF\vw' \to \Coend{\vw' \injto \vw'' \in \commaC{\vw'}\vU} \nat \times \heapF\vw'' \\
    &\countFun(\rho : \va : \vC \injto \vw')(\vh : \heapF\vw') = \quot{1_{\vw'}}\tup\vn\vh \\
    \text{where}& \\
                &  \mathit{go} = \mathsf{fix}\;\Lam{\mathit{rec}} \Lam{\tup{\va \in \vw'}{\vS \subseteq \vw'}} 
                                                                \CaseVertical{\proj\va(\vh)}
                                                                             {\InjL\Unit}{1}
                                                                             {\InjR\vb}{\begin{cases}1 & \text{if } \vb \in \vS \\
                                                                                                     1 + \mathit{rec}(\vb, \{ \vb \} \cup \vS)  & \text{otherwise}\end{cases}} \\
                &  \vn = \mathit{go}\tup\va{\{\va\}}
\end{array}\]
Intuitively, \countFun traverses the heap $\vh$ starting from $\va$ and counts the number of different memory cells that it has encountered.
Notice that \countFun is defined for every extension of $\va : \vC$ to $\vw'$ and so there cannot be a finite $\downset\vn$ that defines it and thus the mapping is not surjective.
\end{proof}

\paragraph{A note about the End in the FGLS monad}

The \countFun function above is well-defined, \ie it satisfies the End relation.
To see why this is the case, notice that the End quantifies over all injections $\rho : \vw \to \vw'$, and thus it gives an $\rho$-ary product: 
\[ \left\{\vf_\rho : \heapF\vw' \to \vP(\Delta_\nat \times \heapF)\vw'\right\}_{\rho : \vw \to \vw' \in \commaC\vw\vU} \]
such that the following coherence condition holds:

For every initialisation of the form 
 \begin{tikzcd} & {\vw} \arrow[dl, "\rho_1"'] \arrow[dr, "\rho_2"] & \\ 
               {\vw_1} \arrow[rr, "\vU\varepsilon"] && {\vw_2} 
 \end{tikzcd} and heap $\vh : \heapF\vw_1$
we have: \[\vP(\Delta_\nat\times\heapF)(\varepsilon)(\vf_{\rho_1}(\vh)) = \vf_{\rho_2}(\heapF\varepsilon\vh)\]

Intuitively, it says each component can only use the reachable part of a heap starting from locations in $\vw$.
That is, take the following heap $\vh$ where the location $\vb$ is not reachable from locations in $\vw$. 
\begin{center}
\begin{tikzpicture}[node distance=5mm]
    \node[pricell=$\vb$] (n1) {};
    \node[mpricell] (n2) [below=0pt of n1] {};
    \node[mpricell] (n3) [below=0pt of n2] {};
    \node[mpricell] (n4) [below=0pt of n3] {};
    \draw[ptr] (n1.center) to ++(0.5,0) to ++(0,-0.55) to (n2.east);
    \draw[ptr] (n3.center) to ++(-0.5,0) to ++(0,0.55) to (n2.west);
    \draw[ptr] (n4.center) to ++(0.5,0) to ++(0,0.55) to (n3.east);
    \draw[brace] (n3.north) to node[labbrace,yshift=-0.1em]{$\vw$} (n4.south);
    \draw[mbrace] (n1.north) to ++\south to node[lambrace,yshift=-0.75em]{$\vw'$} ++\south to++\south to ++\south;
\end{tikzpicture}
\end{center}

The End relation says that the component $\vf$ at world $\vw'$ and heap $\vh$ is completely determined by $\vf$ at $\vw' \setminus \{ \vb \}$ together with $\vh$ without cell $\vb$.
It is then clear to see that \countFun satisfies the End relation as it traverses the heap \emph{starting} from location $\va$.

